\documentclass[12pt,a4paper]{article}

% ============================================================================

\usepackage[utf8]{inputenc}
\usepackage[T5]{fontenc}
\usepackage[vietnamese]{babel}
\usepackage[a4paper,margin=1.8cm]{geometry}
\usepackage{tabularx}
\usepackage{booktabs}
\usepackage{array}
\usepackage{hyperref}
\usepackage{enumitem}
\usepackage{xcolor}
\usepackage{amsmath}
\usepackage{amssymb}
\usepackage{graphicx}
\usepackage{caption}
\usepackage{float}
\usepackage{algorithm}
\usepackage{algpseudocode}

\hypersetup{
    colorlinks=true,
    linkcolor=blue,
    urlcolor=blue,
    pdftitle={EdgeSlice -- Dan y}
}

\setlist[itemize]{leftmargin=*,noitemsep,topsep=4pt}
\setlist[enumerate]{leftmargin=*,noitemsep,topsep=4pt}
\setlength{\parindent}{0pt}
\setlength{\parskip}{6pt}
\renewcommand{\arraystretch}{1.30}

\newcommand{\src}[1]{{\small\texttt{#1}}}
\newcommand{\measured}{\textcolor{blue!60!black}{\small[đã đo]}}
\newcommand{\todo}[1]{\textcolor{red!70!black}{\textbf{[cần đo: #1]}}}

\title{\textbf{Research Report Outline (V9)}\\
\large SRAM-Resident Operator Scheduling for Fast Edge Inference:\\
Capacity-Constrained Tiling, Fusion, and Data Movement in MLIR}
\date{Version: V9 --- revised objective: latency under a hard SRAM constraint}

\begin{document}
\maketitle

\begin{abstract}
\textbf{Abstract:} This report reframes EdgeSlice around the deployment
objective that directly affects performance: every tiled or fused compute
region must fit a declared on-chip SRAM, and legal schedules are ranked by
off-chip traffic and measured latency. Whole-model peak memory remains a
diagnostic rather than the optimization target. The method combines complete
working-set accounting, capacity-constrained 1D/2D tiling, budgeted fusion,
explicit SRAM promotion, double buffering, and target-specific inner-kernel
lowering. Evaluation reports SRAM capacity violations, DRAM--SRAM bytes,
utilization, stalls, latency, energy, and numerical equivalence.
\end{abstract}

\section*{Revised method}
The optimization objective is minimum measured inference latency subject to a
hard on-chip SRAM-capacity constraint. Candidate operations are selected by
their complete local working set---inputs, weights, other operands, halo,
accumulator, temporaries, and output tile---rather than by the whole graph's
peak liveness point. Legal one- and two-dimensional tiles and fusion boundaries
must fit the usable SRAM budget. The compiler then ranks them by predicted
DRAM--SRAM traffic, vector utilization, and parallelism, and measures a short
list on the target.

Tiling alone does not place data in SRAM. The complete backend must allocate a
target memory space, schedule copy-in and copy-out, optionally overlap transfer
and compute with ping-pong buffers, and lower the inner operation to a
vectorized or target-specific kernel. Until that placement is emitted and
verified, a plan is described as \emph{SRAM-feasible}, not SRAM-resident.

\end{document}
